\scanpage{45}%
% archive sheet number: 44
se, if the index theorem is proved for $X$, then $\deleted{\unsure{D|O}}A_{\unsure{\ell|l|\chi}}(X,\handwritten{\unsure{\xi|\zeta}})=B_{\ell}(X)$ is again
independant of the polarization, and analoguous remark for $C_{\ell}(X,\handwritten{\xi})$.

When speaking about condition $C_{\ell}(X,\handwritten{\xi})$, emphasize at once its stability
properties by products (the proof I suggested works indeed) specialization
(with possib\deleted{\unsure{h}}le change of characteristics), hyperplane or more generally
linear sections. Give an extra proposition for the relations with the
property $A$, via a formal proposition as follows:

\uline{Proposition}\quad Condititions équivalentes sur $X$(varieté \uline{polarisée}):

\[
C_{\ell}(X) \tag*{(i)}
\]

\[
C_{\ell}(Y)\text{ et }A_{\ell}(X\mathbin{\mathrm{x}}X)^{\deleted{\unsure{o|0}}},\quad
\deleted{\text{où l'exposant }{}^{o}} \tag*{(ii)}
\]

(ii bis)\quad $C_{\ell}(Y)$ \deleted{\unsure{e}}et $A_{\ell}(X\mathbin{\mathrm{x}}X)^{\unsure{o|0}}$, où l'exposant $^{o}$ signifie qu'on se borne
à exprimer la condition $A$ pour l'homomorphisme en dimension c\deleted{\unsure{ir}}ritique
$H^{2n-2}\handwritten{\longrightarrow}H^{2n+2}$.

\[
C_{\ell}(Y)\text{ et }A_{\ell}(X\mathbin{\mathrm{x}}Y). \tag*{(iii)}
\]

(iii bis)\quad $C_{\ell}(Y)$ et $A_{\ell}(X\mathbin{\mathrm{x}}Y)^{o}$, où l\deleted{\unsure{b}}'exposant $^{o}$ signifie qu'on se borne
à exprimer la condition $A$ en dimension c\deleted{\unsure{ir}}ritique $H^{(2n-1)-1}\handwritten{\longrightarrow}H^{(2n-1)+1}$.

(iv)\quad \deleted{\unsure{Pour tout $0$|Pour tout $O$}} $C_{\ell}(Y)$, et pour tout $i\handwritten{\leq} n-1$, l'homomorphisme
naturel $H^{i}(Y)\handwritten{\unsure{-|\to}}H^{i}(X)$ inverse à gauche de $\phi^{i}:H^{i}(X)\handwritten{-}H^{i}(Y)$ (induit
par $\handwritten{\Lambda}_{X}\phi_{\handwritten{*}}$) est induit par une classe de correspondance algébrique (induisant ce qu'elle veut sur les autres $H^{j}(Y)$ ).

(iv bis)\quad $C_{\ell}(Y)$, et pour $j\handwritten{\geq} n+1$, l'homomorphisme naturel $H^{j}(X)\handwritten{\to}H^{\unsure{j-2}}(Y$
inverse à droite de $\phi_{j-2}:H^{j-2}(Y)\handwritten{\unsure{-|\to}}H^{j}(X)$ \deleted{est} (induit par $\phi^{*}\handwritten{\Lambda}_{X}$) est
induit par une classe de correspondance algébrique (induisant ce qu'elle
veut sur les autres $H^{i}(X)$).

\uline{Corollaire}\quad \deleted{\unsure{L}}Ces conditions équivalent aussi à

(v)\quad $A_{\ell}(X\mathbin{\mathrm{x}}X)^{\unsure{o|\,|0}}+A_{\handwritten{\ell}}(Y\mathbin{\mathrm{x}}Y)^{\handwritten{o}}+A_{\handwritten{\ell}}(Z\mathbin{\mathrm{x}}Z)^{\handwritten{o}}+\hspace{2em}$, où $X\handwritten{\supset}Y\handwritten{\supset}Z$ est une
suite décorissante de sections hyperplanes\handwritten{\unsure{ non sing.}}

